% =====================================================================
% ch02a —— 第 2 章章首 + §2.1 定义 + §2.2 周期性定理（前段）
% 原书 p.29–p.38（PDF 38–47）
% =====================================================================
% [原文 p.29]
\chapter{K 理论}
\section{定义}
\takeaway{对紧空间上的向量丛半群作 Grothendieck 构造，得到环 $K(X)$；本节给出 $K(A)$ 的两种构造，并说明 $[E]=[F]$ 恰好等价于两丛稳定等价。}

设 $X$ 为紧 Hausdorff 空间，则集合 $\Vect(X)$ 具有交换半群的结构，其加法结构由直和定义。设 $A$ 为任意交换半群，则可以给 $A$ 配上一个具有如下性质的 Abel 群 $K(A)$：存在半群同态 $\alpha\colon A\to K(A)$，使得对任意群 $G$ 与任意半群同态 $\gamma\colon A\to G$，都存在唯一的同态映射 $\chi\colon K(A)\to G$ 满足 $\gamma=\chi\circ\alpha$。这样的 $K(A)$ 若存在，则在差一个同构的意义下唯一。

群 $K(A)$ 按如下方式构造。设 $F(A)$ 为由 $A$ 的元素生成的自由 Abel 群，以 $\oplus$ 表示 $A$ 的加法。设 $E(A)$ 为由对 $a,a'\in A$ 而言形如 $a+a'-(a\oplus a')$ 的元素所生成的 $F(A)$ 的子群。此时 $K(A)=F(A)/E(A)$，取 $\alpha\colon A\to K(A)$ 为商映射，就具有前述的普遍性。

还有与之稍异、但在某些场合方便的 $K(A)$ 的如下构造法。设 $\Delta\colon A\to A\times A$ 为对角映射，它是半群同态。以 $K(A)$ 表示 $A\times A$ 中 $\Delta(A)$ 的剩余类的集合。这是商半群。然而，交换 $A\times A$ 的两个分量便得到逆元，所以 $K(A)$ 是群。此时，把 $\alpha_A\colon A\to K(A)$ 定义为 $a\mapsto(a,0)$ 与自然投影 $A\times A\to K(A)$ 的合成（为简单起见，假定 $A$ 有零元）。对 $(K(A),\alpha_A)$ 关于 $A$ 是函子式的，因此当 $\gamma\colon A\to B$ 为半群同态时，图式
% [原文 p.30]
\[
\begin{array}{ccc}
A & \xrightarrow{\ \alpha_A\ } & K(A)\\
{\scriptstyle\gamma}\downarrow & & \downarrow{\scriptstyle K(\gamma)}\\
B & \xrightarrow{\ \alpha_B\ } & K(B)
\end{array}
\]
是可交换的。若 $B$ 是群，则 $\alpha_B$ 是同构映射。这表明 $K(A)$ 具有所求的普遍性。

若 $A$ 还是半环（即 $A$ 具有对加法分配的乘法），则 $K(A)$ 显然是环。

设 $X$ 为空间。记环 $K(\Vect(X))$ 为 $K(X)$。这个记号大概不会引起混乱。设 $E\in\Vect(X)$，把 $E$ 在 $K(X)$ 中的像记作 $[E]$。为了避免记号繁杂，在没有混淆之虞时，也把 $[E]$ 简写成 $E$。

% [Tutor Note: 普遍性只保证 $\alpha\colon A\to K(A)$ 是半群同态，并不保证它是单射；$A$ 若没有消去律，不同的元素可能在 $K(A)$ 中有同一个像。对 $A=\Vect(X)$ 也如此，这正是下文必须引入稳定等价的原因。]
用 $K$ 群的第二种构造法，设 $X$ 为空间，则 $K(X)$ 的所有元素都可以表成形如 $[E]-[F]$ 的元素，其中 $E$，$F$ 是 $X$ 上的向量丛。设 $G$ 是使 $F\oplus G$ 平凡的向量丛。把 $n$ 维平凡向量丛记作 $\underline{n}$。设 $F\oplus G=\underline{n}$，则
\[
[E]-[F]=[E]+[G]-([F]+[G])=[E\oplus G]-[\underline{n}].
\]
这样，$K(X)$ 的所有元素也可以表成形如 $[H]-[\underline{n}]$ 的元素。

% [Tutor Note: 带下划线的 $\underline{n}$ 专指 $n$ 维平凡向量丛，与整数 $n$ 区分；后文还会出现 $[\underline{n}]$ 与 $\underline{n}$ 混用之处，读者须按上下文判断哪个是 $K(X)$ 中的类。]
若 $E$，$F$ 满足 $[E]=[F]$，则由 $K$ 群的第二种构造法再次可知，存在向量丛 $G$ 使 $E\oplus G\cong F\oplus G$。设 $G'$ 是使 $G\oplus G'\cong\underline{n}$ 的向量丛。此时 $E\oplus G\oplus G'\cong F\oplus G\oplus G'$，故 $E\oplus\underline{n}\cong F\oplus\underline{n}$。当两个向量丛各自加上适当的平凡丛后同构时，称这两个向量丛稳定等价（stably equivalent）。因此，$[E]=[F]$ 当且仅当 $E$ 与 $F$ 稳定等价。

设 $f\colon X\to Y$ 为连续映射。此时 $f^{*}\colon\Vect(Y)\to\Vect(X)$ 诱导环同态 $f^{*}\colon K(Y)\to K(X)$。由引理 1.4.3，这个环同态只依赖于 $f$ 的同伦类。

% [原文 p.31]
\section{周期性定理}
\takeaway{周期性定理断言 $K(X)\otimes K(S^2)\cong K(X\times S^2)$。本节引入伴随射影丛、零截面与无穷远截面、以及 Laurent 衔接函数，把定理的证明归约为对衔接函数形状的分析。}

$K$ 理论的基本定理是周期性定理。它最简单的形式断言：对任意紧 Hausdorff 空间 $X$，$K(X)\otimes K(S^2)$ 与 $K(X\times S^2)$ 之间存在同构映射。这是后面要证明的更一般定理的特别情形。

设 $E$ 是空间 $X$ 上的向量丛，把 $X$ 看作 $E$ 中的零截面，令 $E_0=E\setminus X$。此时，非零复数作为保纤维的自同构映射作用于 $E_0$。用这个作用作商，从 $E_0$ 得到商空间 $P(E)$。$P(E)$ 称为 $E$ 的伴随射影丛。设 $p\colon P(E)\to X$ 为投影，则对每个 $x\in X$，$p^{-1}(x)$ 是复射影空间。

设 $V$ 为向量空间，$W$ 为一维向量空间，则 $V$ 与 $V\otimes W$ 同构，但并没有自然的同构。对任意非零元素 $w\in W$，映射 $v\mapsto v\otimes w$ 定义了 $V$ 与 $V\otimes W$ 之间的同构映射，因此定义了同构映射 $P(w)\colon P(V)\to P(V\otimes W)$。又设 $w'$ 是 $W$ 的另一个非零元素，则存在非零复数 $\lambda$ 使 $w'=\lambda w$。因此 $P(w)=P(w')$，从而 $P(V)$ 与 $P(V\otimes W)$ 之间的同构映射是自然的。所以，若 $E$ 是任意向量丛、$L$ 是线丛，则存在自然同构映射 $P(E)\cong P(E\otimes L)$。

设 $E$ 为 $X$ 上的向量丛，则每一点 $a\in P(E)_x=P(E_x)$ 代表一个一维子空间 $H_a^{*}\subset E_x$。设 $p\colon P(E)\to X$ 为投影，则这些子空间全体的并集是子空间 $H^{*}\subset p^{*}E$。$H^{*}$ 是 $p^{*}E$ 的子丛，这一点容易验证。事实上，问题局部化后可以假定 $E$ 是积丛，于是归结为 1.4 节已经讨论过的 Grassmann 流形的特别情形。把这个线丛记作 $H^{*}$，是因为想要它的对偶 $H$（这里记号的选取由本书不讨论的代数几何方面的考虑决定）。

于是可以叙述周期性定理了。
\begin{theorem*}{2.2.1}
  设 $L$ 为 $X$ 上的线丛。此时作为 $K(X)$ 代数，
% [原文 p.32]
  $K(P(L\oplus\underline{1}))$ 由 $[H]$ 生成，且只服从唯一的关系式
  \[
  ([H]-[\underline{1}])([L][H]-[\underline{1}])=0.
  \]
\end{theorem*}

在进入这个定理的证明之前，先举出两个推论。注意 $P(\underline{1}\oplus\underline{1})=X\times S^2$。

\begin{corollary*}{2.2.2}
  $K(S^2)$ 作为 $K(*)$ 模由 $[H]$ 生成，而 $[H]$ 只服从唯一的关系式 $([H]-[\underline{1}])^2=0$。
\end{corollary*}

\begin{corollary*}{2.2.3}
  对任意空间 $X$，当 $\mu\colon K(X)\otimes K(S^2)\to K(X\times S^2)$ 由 $\mu(a\otimes b)=(\pi_1^{*}a)(\pi_2^{*}b)$ 定义时，$\mu$ 是环同构映射。其中 $\pi_1$，$\pi_2$ 是两个因子上的投影。
\end{corollary*}

定理 2.2.1 的证明分解为一连串引理。

首先，对任意 $x\in X$，注意自然嵌入 $L_x\to P(L\oplus\underline{1})_x$ 由映射 $y\mapsto(y,1)$ 给出。这个映射可延拓到 $L_x$ 的单点紧化，给出从 $L_x$ 的单点紧化到 $P(L\oplus\underline{1})_x$ 上的同胚映射。通过把 $x$ 映到 $L_x$ 的单点紧化的「无穷远点」的像来定义映射 $X\to P(L\oplus\underline{1})$，便得到 $P(L\oplus\underline{1})$ 的一个截面，称为「无穷远截面」。同样地，$L$ 的零截面也给出 $P(L\oplus\underline{1})$ 的一个截面，称之为 $P(L\oplus\underline{1})$ 的零截面。

给 $L$ 配上度量，设 $S\subset L$ 为单位圆周丛。把 $L$ 中长度不超过 1 的向量组成的子集记作 $P^0$；把由无穷远截面以及长度不小于 1 的所有向量组成的 $P(L\oplus\underline{1})$ 的子集记作 $P^\infty$。射影 $S\to X$，$P^0\to X$，$P^\infty\to X$ 分别记作 $\pi$，$\pi_0$，$\pi_\infty$。

% [Tutor Note: 这里的 $P^0$ 是单位圆盘丛（长度不超过 1 的部分），$P^\infty$ 是补上无穷远截面后的外部，两者都是 $P(L\oplus\underline{1})$ 的闭子集；$S$ 是它们的公共边界。零截面落在 $P^0$ 内，无穷远截面是 $P^\infty$ 的「尖端」。]
由于 $\pi_0$ 与 $\pi_\infty$ 同伦等价，$P^0$ 上的所有向量丛都呈 $\pi_0^{*}(E^0)$ 的形状，$P^\infty$ 上的所有向量丛都呈 $\pi_\infty^{*}(E^\infty)$ 的形状，其中 $E^0$ 与 $E^\infty$ 是 $X$ 上的向量丛。因此，$P(L\oplus\underline{1})$ 上的任意向量丛 $E$ 都同构于以 $f\in\mathrm{ISO}(\pi^{*}(E^0),\pi^{*}(E^\infty))$ 为衔接函数的形如 $(\pi_0^{*}(E^0),f,\pi_\infty^{*}(E^\infty))$ 的向量丛。此外，同构映射
\[
E\longrightarrow(\pi_0^{*}(E^0),f,\pi_\infty^{*}(E^\infty))
\]
% [原文 p.33]
在零截面与无穷远截面上都与自明映射一致，则可推出 $f$ 的同伦类由 $E$ 的同构类唯一决定。为此，再次用到零截面是 $P^0$ 的形变收缩、无穷远截面是 $P^\infty$ 的形变收缩这两个事实。以下把 $(\pi_0^{*}(E^0),f,\pi_\infty^{*}(E^\infty))$ 简记作 $(E^0,f,E^\infty)$，使记号稍微简单些。

现在要证明的是：向量丛 $E^0$ 与 $E^\infty$，以及衔接函数 $f$，都可以取成极为简单的形状。在 $L$ 平凡的特别情形，$S$ 就是 $X\times S^1$，而投影 $S\to S^1$ 是 $S$ 上的复值函数。把这个函数记作 $z$（这里把 $S^1$ 等同于绝对值等于 1 的复数集合）。于是，可以以 $X$ 上的函数为系数，考虑作为 $z$ 的有限 Laurent 级数的 $S$ 上的函数
\[
\sum_{k=-n}^{n}a_k(x)z^k.
\]
这些有限 Laurent 级数可以一致逼近 $S$ 上的函数。

当 $L$ 非平凡时，也有有限 Laurent 级数的类似物。此时 $z$ 不再是函数，而是向量丛的截面。$\pi^{*}(L)$ 是 $S\times L$ 的部分集合，故对角映射 $S\to S\times S\subset S\times L$ 给出 $\pi^{*}(L)$ 的截面。把这个截面记作 $z$。作张量积，则对 $k\ge0$ 得到 $(\pi^{*}(L))^k$ 的截面 $z^k$，以及 $(\pi^{*}(L^{*}))^k$ 的截面 $z^{-k}$。把 $(L^{*})^k$ 记作 $L^{-k}$，则对任意整数 $k,k'$ 有 $L^k\otimes L^{k'}\cong L^{k+k'}$。设 $a_k\in L^{-k}$，则 $\pi^{*}(a_k)\otimes z^k\in\Gamma(\pi^{*}(\underline{1}))$，因此 $\pi^{*}(a_k)\otimes z^k$ 是 $S$ 上的函数。把这个函数记作 $a_kz^k$。$S$ 上的有限 Laurent 级数，是指形如
\[
\sum_{k=-n}^{n}a_kz^k
\]
的 $S$ 上函数之和；其中对所有 $k$ 都有 $a_k\in\Gamma(L^{-k})$。

% [Tutor Note: 原书此句前半作 $a_k\in L^{-k}$，紧接的后一句却写作 $a_k\in\Gamma(L^{-k})$；按上下文，$a_k$ 是 $L^{-k}$ 的截面，前一处漏排了 $\Gamma$。译文照录原文，读者按截面理解即可。]
更一般地，当 $E^0$，$E^\infty$ 是 $X$ 上的向量丛，且 $a_k\in\Gamma(\Hom(L^k\otimes E^0,E^\infty))$ 时，把 $a_k\otimes z^k$ 记作 $a_kz^k$，则
% [原文 p.34]
形如
\[
f=\sum_{k=-n}^{n}a_kz^k
\]
的有限和是 $\Gamma(\pi^{*}(E^0),\pi^{*}(E^\infty))$ 的元素。更进一步，当 $f\in\mathrm{ISO}(\pi^{*}(E^0),\pi^{*}(E^\infty))$ 时，称 $f$ 为关于 $(E^0,E^\infty)$ 的 Laurent 衔接函数（Laurent clutching function）。

% [Tutor Note: 原书把 Hom 空间的截面集记作 $\Gamma(\pi^{*}(E^0),\pi^{*}(E^\infty))$，但 $\Gamma$ 只取一个变元，这个记号与 p.20 的 $\mathrm{HOM}$ 用法不一致；按紧接的 $f\in\mathrm{ISO}(\cdots)$ 判断，此处指的是丛同态 $\Hom(\pi^{*}(E^0),\pi^{*}(E^\infty))$ 的截面。]
函数 $z$ 是关于 $(\underline{1},L)$ 的（Laurent）衔接函数。更进一步，$(\underline{1},z,L)$ 正是已经定义过的向量丛 $H^*$。要看出这一点，首先回忆 $H^*$ 是作为 $\pi^*(L\oplus\underline{1})$ 的子丛定义的。对每一点 $y\in P(L\oplus\underline{1})_x$，$H^*_y$ 是 $(L\oplus\underline{1})_x$ 的子空间，且
\[
H^*_\infty=L_x\oplus0,\qquad H^*_0=0\oplus\underline{1}_x
\]
成立。因此，合成
\[
H^*\longrightarrow\pi^*(L\oplus\underline{1})\longrightarrow\pi^*(\underline{1})
\]
由射影 $L\oplus\underline{1}\to\underline{1}$ 诱导，并确定同构映射
\[
f_0\colon H^*|_{P^0}\longrightarrow\pi_0^*(\underline{1}).
\]
同样地，合成
\[
H^*\longrightarrow\pi^*(L\oplus\underline{1})\longrightarrow\pi^*(L)
\]
由射影 $L\oplus\underline{1}\to L$ 诱导，并确定同构映射
\[
f_\infty\colon H^*|_{P^\infty}\longrightarrow\pi_0^*(L).
\]
% [Tutor Note: 原书 p.34 把 $f_\infty$ 的靶印作 $\pi_0^*(L)$，但 $\pi_0$ 是 $P^0$ 上的投影，$f_\infty$ 定义在 $P^\infty$ 上，此处应为 $\pi_\infty^*(L)$，疑为笔误。译文照录原文。]
此时，$f=f_\infty\circ f_0^{-1}\colon\pi^*(\underline{1})\to\pi^*(L)$ 是关于 $H^*$ 的衔接函数。显然，对 $y\in S_x$，$f(y)$ 是同构映射，其图像是 $H^*_y$。$H^*_y$ 是由 $y\oplus1$（$y\in S_x\subset L_x$，$1\in\C$）张成的 $L_x\oplus\underline{1}_x$ 的子空间，因此可知 $f$ 就是这个截面 $z$ 本身。所以
\[
H^*\cong(\underline{1},z,L)
\]
成立。因此，对任意整数 $k$，
% [原文 p.35]
\[
H^k\cong(\underline{1},z^{-k},L^{-k})
\]
成立。

把 $P=P(L\oplus\underline{1})$ 上的向量丛分类的下一阶段，是证明所有的衔接函数都可以取成 Laurent 衔接函数。设 $f\in\Gamma(\Hom(\pi^{*}E^0,\pi^{*}E^\infty))$ 为任意截面。$f$ 的 Fourier 系数由
\[
a_k\in\Gamma(\Hom(L^k\otimes E^0,E^\infty)),
\]
\[
a_k(x)=\frac{1}{2\pi i}\int_{S_x}f_xz_x^{-k-1}dz_x
\]
定义。这里，$f_x$，$z_x$ 表示 $f$，$z$ 在 $S_x$ 上的限制，因此 $dz_x$ 是 $S_x$ 上的、以 $L_x$ 为系数的微分形式。设 $S_n$ 为部分和
\[
S_n=\sum_{k=-n}^{n}a_kz^k,
\]
并定义 Cesàro 平均为
\[
f_n=\frac{1}{n}\sum_{k=0}^{n-1}S_k.
\]
此时，关于 Fourier 级数的 Cesàro 平均的 Fejér 定理的证明可以立即推广到这里的一般情形，从而得到下述引理。

% [Tutor Note: 这里 $f$ 是紧空间 $S$ 上的连续截面，Fourier 系数用逐纤维的围道积分定义；系数 $a_k$ 落在 $\Gamma(\Hom(L^k\otimes E^0,E^\infty))$ 中，是因为 $f_xz_x^{-k-1}dz_x$ 沿 $S_x$ 积分后在 $L_x^{-k}$ 中取值，再与 $E^0_x$ 张量并作用到 $E^\infty_x$。]
\begin{lemma*}{2.2.4}
  设 $f$ 为关于 $(E^0,E^\infty)$ 的衔接函数，$f_n$ 为 $f$ 的 Fourier 级数的 Cesàro 平均序列。此时 $f_n$ 一致收敛于 $f$。因此，对充分大的所有 $n$，$f_n$ 是关于 $(E^0,E^\infty)$ 的衔接函数，且 $(E^0,f,E^\infty)\cong(E^0,f_n,E^\infty)$ 成立。
\end{lemma*}

\begin{proof}
  由于 $\mathrm{ISO}(E^0,E^\infty)$ 是向量空间 $\mathrm{HOM}(E^0,E^\infty)$ 的开集，存在 $\varepsilon>0$，使得当 $|f-g|<\varepsilon$ 时 $g\in\mathrm{ISO}(E^0,E^\infty)$。其中 $|\cdot|$ 是关于 $E^0$，$E^\infty$ 的固定度量的一致范数（上确界范数）。由于 $f_n$ 一致收敛于 $f$，对充分大的 $n$ 有 $|f-f_n|<\varepsilon$。因此，对 $0\le t\le1$ 有 $tf+(1-t)f_n\in\mathrm{ISO}(E^0,E^\infty)$。所以 $f$ 与 $f_n$ 在 $\mathrm{ISO}(E^0,E^\infty)$ 中同伦，故 $(E^0,f,E^\infty)\cong(E^0,f_n,E^\infty)$
% [原文 p.36]
成立。
\end{proof}

接下来，考虑多项式衔接函数
\[
p=\sum_{k=0}^{n}a_kz^k.
\]
考虑由矩阵
\[
\mathcal{L}^n(p)=
\begin{pmatrix}
a_0 & a_1 & a_2 & \cdots & a_n\\
-z & 1 & & & \\
& -z & 1 & & \\
& & \ddots & \ddots & \\
& & & -z & 1
\end{pmatrix}
\]
给出的同态映射
\[
\mathcal{L}^n(p)\colon\pi^{*}\left(\sum_{k=0}^{n}L^k\otimes E^0\right)\longrightarrow\pi^{*}\left(E^\infty\oplus\sum_{k=1}^{n}L^k\otimes E^0\right).
\]

显然，$\mathcal{L}^n(p)$ 关于 $z$ 是线性的。用
\[
p_0=p,\qquad zp_{r+1}(z)=p_r(z)-p_r(0)
\]
递归地定义列 $p_r(z)$。于是得到如下矩阵恒等式：
\[
\mathcal{L}^n(p)=
\begin{pmatrix}
1 & p_1 & p_2 & \cdots & p_n\\
& 1 & & & \\
& & 1 & & \\
& & & \ddots & \\
& & & & 1
\end{pmatrix}
\begin{pmatrix}
p & & & & \\
& 1 & & & \\
& & 1 & & \\
& & & \ddots & \\
& & & & 1
\end{pmatrix}
\begin{pmatrix}
1 & & & & \\
-z & 1 & & & \\
& -z & 1 & & \\
& & \ddots & \ddots & \\
& & & -z & 1
\end{pmatrix}.
\]
把它写得更简单些，就是
\[
\mathcal{L}^n(p)=(1+N_1)(p\oplus1)(1+N_2),
\]
其中 $N_1$ 与 $N_2$ 是幂零的。若 $N$ 幂零，则对 $0\le t\le1$，$1+tN$ 非奇异，因此得到下述命题。
% [原文 p.37]
\begin{proposition*}{2.2.5}
  $\mathcal{L}^n(p)$ 与 $p\oplus1$ 定义了 $P$ 上同构的向量丛。即
  \[
  \begin{gathered}
  (E^0,p,E^\infty)\oplus\left(\sum_{k=1}^{n}L^k\otimes E^0,1,\sum_{k=1}^{n}L^k\otimes E^0\right)\\
  \cong\left(\sum_{k=0}^{n}L^k\otimes E^0,\mathcal{L}^n(p),E^\infty\oplus\sum_{k=1}^{n}L^k\otimes E^0\right)
  \end{gathered}
  \]
  成立。
\end{proposition*}

\begin{remark}{}
  当然，$\mathcal{L}^n(p)$ 的定义模仿了把 $n$ 阶常微分方程化为一阶方程组的方法。
\end{remark}

为简单起见，把向量丛
\[
\left(\sum_{k=0}^{n}L^k\otimes E^0,\mathcal{L}^n(p),E^\infty\oplus\sum_{k=1}^{n}L^k\otimes E^0\right)
\]
记作 $\mathcal{L}^n(E^0,p,E^\infty)$。

\begin{lemma*}{2.2.6}
  设 $p$ 为关于 $(E^0,E^\infty)$ 的次数不超过 $n$ 的多项式衔接函数。此时
  \begin{enumerate}[label=(\roman*),leftmargin=2.2em,itemsep=2pt]
  \item $\mathcal{L}^{n+1}(E^0,p,E^\infty)\cong\mathcal{L}^n(E^0,p,E^\infty)\oplus(L^{n+1}\otimes E^0,1,L^{n+1}\otimes E^0)$；
  \item $\mathcal{L}^{n+1}(L^{-1}\otimes E^0,zp,E^\infty)\cong\mathcal{L}^n(E^0,p,E^\infty)\oplus(L^{-1}\otimes E^0,z,E^0)$
  \end{enumerate}
  成立。
\end{lemma*}

\begin{proof}
  首先，
  \[
  \mathcal{L}^{n+1}(p)=
  \begin{pmatrix}
  \mathcal{L}^n(p) & 0\\
  0\ \ 0\ \ \cdots\ \ -z & 1
  \end{pmatrix}
  \]
  成立。把最下行的 $z$ 乘以 $t$，便得到 $\mathcal{L}^{n+1}(p)$ 与 $\mathcal{L}^n(p)\oplus1$ 之间的同伦。这样 (i) 得证。

  同样地，
% [原文 p.38]
\[
\mathcal{L}^{n+1}(zp)=
\begin{pmatrix}
0 & a_0 & a_1 & \cdots & a_n\\
-z & 1 & & & \\
& -z & 1 & & \\
& & \ddots & \ddots & \\
& & & -z & 1
\end{pmatrix}
\]
成立。把这个第 2 行的 $1$ 乘以 $t$，便得到 $\mathcal{L}^{n+1}(zp)$ 与 $\mathcal{L}^n(p)\oplus(-z)$ 之间的同伦。$-z$ 是 $z$ 与映射 $-1$ 的合成，而 $-1$ 可延拓到 $E^0$，所以 $(L^{-1}\otimes E^0,-z,E^0)\cong(L^{-1}\otimes E^0,z,E^0)$。因此 (ii) 得证。
\end{proof}

下面证明 $K(P)$ 中的一个简单关系式。把 $[(E^0,p,E^\infty)]$ 记作 $[E^0,p,E^\infty]$。

\begin{proposition*}{2.2.7}
  对关于 $(E^0,E^\infty)$ 的任意多项式衔接函数 $p$，等式
  \[
  ([E^0,p,E^\infty]-[E^0,1,E^0])([L][H]-[\underline{1}])=0
  \]
  成立。
\end{proposition*}

\begin{proof}
  由引理 2.2.6 的后半部分与命题 2.2.7 合起来，可知
% [Tutor Note: 原书 p.38 此处作「命題 2.2.7」，但这一段正在证明命题 2.2.7，自引不合逻辑；由所引内容（$\mathcal{L}^n(p)$ 与 $p\oplus1$ 界定同构的丛）判断，应为命题 2.2.5。译文照录原文。]
  \[
  \begin{gathered}
  (L^{-1}\otimes E^0,zp,E^\infty)\oplus\left(\sum_{k=0}^{n}L^k\otimes E^0,1,\sum_{k=0}^{n}L^k\otimes E^0\right)\\
  \cong(E^0,p,E^\infty)\oplus\left(\sum_{k=1}^{n}L^k\otimes E^0,1,\sum_{k=1}^{n}L^k\otimes E^0\right)\oplus(L^{-1}\otimes E^0,z,E^0)
  \end{gathered}
  \]
  成立。因此，在 $K(P)$ 中，
  \[
  [L^{-1}\otimes E^0,zp,E^\infty]\oplus[E^0,1,E^0]=[E^0,p,E^\infty]\oplus[L^{-1}\otimes E^0,z,E^0]
  \]
  成立。由于 $[\underline{1},z,L]=[H^{-1}]$，
  \begin{equation*}
  \resizebox{0.96\textwidth}{!}{$\displaystyle [L^{-1}][H^{-1}][E^0,p,E^\infty]\oplus[E^0,1,E^0]=[E^0,p,E^\infty]\oplus[L^{-1}][H^{-1}][E^0,1,E^0]$}
  \end{equation*}
  成立。
\end{proof}

% 块尾跨页句：「[L^{-1}][H^{-1}][E^0,p,E^\infty]\oplus[E^0,1,E^0]=...」起于 p.38，句末「成立。」与 □ 在 p.39 顶部；本块已收尾（由 ch02b 续写 p.39 正文）。
% CHUNK ch02a DONE